\section{Introduction}
\label{sec:intro}

Diffusion models generate by integrating a reverse noising process whose only learned component is a score network $s_\theta(x_t,t)\approx\nabla\log p_t(x_t)$~\citep{ho2020denoising,song2021scoresde,karras2022elucidating}. Latent diffusion runs this process in a learned latent space of dimension $\dlat$~\citep{rombach2022highresolution,esser2024sd3} that still exceeds the intrinsic dimension $\dint$ of the data~\citep{stanczuk2024manifold,kamkari2024geometric}. A network that fits the empirical score of a finite training set reproduces that set; practical models generalize only for a while, and recent work describes training as a transition from population learning to sample collapse~\citep{yi2023generalization,biroli2024dynamical,bonnaire2025}. Memorization is therefore both a privacy and copyright problem~\citep{carlini2019secret,somepalli2023diffusion,carlini2023extracting} and a basic question about what the model learns. The latent dimension sits between the two, and its effect on when memorization begins has not been isolated.

We ask whether widening the latent space, with everything else fixed, delays memorization, and when a spectral model can explain the delay. On the first part the evidence is clear. For MLP score networks of fixed hidden width trained on synthetic mixtures with $\dint=5$, the fraction of memorized samples after 5M steps falls from $30\%$ at $\dlat=5$ to near zero ($0.1\%$) at $\dlat=40$ over five seeds, and the learned score departs from the population score progressively later as the latent widens (\cref{sec:phenomenon}). The decline is not the nearest-neighbor test losing sensitivity in high dimension: restricted to the true five-dimensional signal subspace, where the signal-copy control retains its sensitivity, memorization still falls from $29.9\%$ to $1.2\%$ against a chance level of $0.4\%$ (\cref{tab:projection-main}). On CelebA and CIFAR-10 VAE latents, once the VAE is wide enough to preserve sample identity, the mean time to reach $5\%$ memorization increases by $3.6\times$ and $1.4\times$, respectively, over post-threshold width comparisons, while FID follows a separate quality tradeoff (\cref{sec:real}). On the second part the answer is partial: a solvable frozen-feature model identifies one mechanism for the delay (\cref{sec:theory}), but its transfer to learned representations is not established (\cref{sec:real,sec:discussion}).

Our contributions are:
\begin{itemize}
  \item \textbf{The phenomenon, with a dimension-controlled test.} A fixed-width $\dlat$ sweep at known $\dint$ in which memorization falls from $30\%$ to near zero, together with a projected variant of the nearest-neighbor ratio test. The variant is necessary: a sampler that copies signal coordinates and redraws null coordinates is invisible to the full-space test for $\dlat\ge10$, so that test alone cannot separate a real delay from a loss of sensitivity (\cref{sec:setup,sec:phenomenon}).
  \item \textbf{A spectral explanation in frozen features.} Evaluating the random-feature theory of \citet{bonnaire2025} at the two-atom covariance of data embedded in a larger latent, the feature-correlation spectrum splits into signal, noise-dimension, sample-specific and null blocks (\cref{prop:fourblock-9}). Under gradient flow, modes learn in parallel and the measured ratio of signal to sample-block edges grows with $\dlat$ as the pre-activation variance $q$ falls and nonlinear feature mass shrinks. Interpreting this ratio as a memorization clock requires spectral-alignment assumptions. A control that holds $q$ fixed as $\dlat$ grows largely removes the endpoint growth ($1.09\times$ against $3.17\times$ over $\dlat=40\to240$), supporting the operating-point explanation (\cref{sec:theory}).
  \item \textbf{Real-data evidence and the transfer boundary.} CelebA and CIFAR-10 sweeps quantify memorization timing and image quality after identity preservation. Measurements in trained MLPs and a spatial-DiT pilot delimit the transfer of the frozen-feature explanation (\cref{sec:real}).
\end{itemize}

\paragraph{Relation to prior work.}
\citet{bonnaire2025} characterize diffusion training in a random-feature network as a competition between a generalization time $\taugen$ and a memorization time $\taumem$. Their Theorem~3.1 gives the limiting feature spectrum for an arbitrary input spectral measure, but their experiments and two-bulk reading take the isotropic case, which collapses the distinction between the $\dint$ signal directions and the $\dlat-\dint$ null directions. \citet{george2025dsm} analyze a linear pencil for data supported exactly on a $\dint$-dimensional subspace, the $\signoise=0$ limit. We extend neither result: we evaluate the existing theorem at the two-atom covariance with strictly positive null variance $\signoise>0$ and read the resulting blocks as gradient-flow learning-rate bands. Memorization in generative models is otherwise studied through membership inference, reconstruction attacks and data-copying tests~\citep{hayes2017logan,hilprecht2019reconstruction,meehan2020datacopying}, and for diffusion models through nearest-neighbor replication and extraction~\citep{somepalli2023diffusion,carlini2023extracting}; we use the nearest-neighbor ratio criterion in the form of \citet{bonnaire2025}.
